% Source: ACM_Templates.pdf, PDF pp. 7-8, string hash/wxf.
% Validation and changes: docs/validation.md; scripts/check_strings.py.
\subsection{字符串双哈希}
\label{algo:string-hash}
\AlgorithmInfo
    {快速比较子串，支持在一次预处理后反复查询任意区间的哈希。}
    {预处理时间、空间均为 $O(n)$；每次区间查询或比较 $O(1)$。}
    {传入原字符串，查询为从 1 开始的闭区间；空区间返回 $(0,0)$。}

\subsubsection*{原理}
将字符串看作以 $B$ 为底的多项式，维护前缀哈希 $h_i$ 和幂 $p_i=B^i$。设字符编码为 $c_i$，则
\[
    h_i=(h_{i-1}B+c_i)\bmod M,\qquad
    H(l,r)=(h_r-h_{l-1}p_{r-l+1})\bmod M.
\]
减去前面的贡献后，相同子串无论位于哪里都会得到相同结果。模板使用两组底数、模数，把两个结果一起比较以降低碰撞概率。

\subsubsection*{使用教程}
\InlineCode{init(s)} 建立前缀；\InlineCode{query(l,r)} 返回闭区间的双哈希数对，\InlineCode{same(l1,r1,l2,r2)} 返回同一字符串内两个子串是否等长且哈希相同。
\begin{lstlisting}[style=mtycpp,numbers=none]
StringHash ha;
ha.init("ababa");
auto h=ha.query(1,3);
bool ok=ha.same(1,3,3,5);
\end{lstlisting}
比较不同字符串时，分别初始化两个对象，在确认区间长度相同后比较两次 \InlineCode{query} 的结果。\InlineCode{init} 可重复调用。

\needspace{7\baselineskip}
\subsubsection*{模板代码}
\CodeFile{code/strings/hash.cpp}

\subsubsection*{需要注意}
哈希相同仍存在碰撞可能，双哈希也不提供绝对保证；需要确定性匹配时使用 KMP 等算法。非空区间须满足 $1\leq l\leq r\leq n$；不要把普通字符串的从 0 开始下标直接传给查询接口。
